\paragraph{平行族截断塔上的相干停时、信息增益与 \(4\)-进读出}\label{par:xi-terminal-zm-leyang-coherence-stopping-time}
本段把 Lee--Yang 分歧塔的“平行族并/截断语义”抽象为一个 \(4\)-叉纤维的有限群塔接口，并给出相干停时的精确律、分辨率信息恒等式以及与 \(4\)-进赋值的严格同构。

\begin{theorem}[平行族 \(4\)-叉群塔上的二元相干停时律]\label{thm:xi-coherence-stopping-time-k-parallel-4ary}
\leanverified{paper\_xi\_coherence\_stopping\_time\_k\_parallel\_4ary}
设 \((G_n)_{n\ge 1}\) 为有限阿贝尔群塔，满足 \(|G_n|=4^n\)。
对每个 \(n\ge 1\) 给定满射
\(
\pi_{n+1\to n}:G_{n+1}\twoheadrightarrow G_n
\)
且每个纤维基数恒为 \(4\)。
固定整数 \(k\ge 1\)，并对每个 \(n\ge 1\) 取集合
\[
A_n:=\bigsqcup_{i=1}^{k}A_{n,i},
\qquad
A_{n,i}\cong G_n,
\]
其中每个 \(A_{n,i}\) 是一个 \(G_n\)-挠空间（主齐次空间）。
假设存在兼容的截断族 \(\Pi_{n+1\to n}:A_{n+1}\to A_n\)，使得对每个 \(i\)：
\[
\Pi_{n+1\to n}(A_{n+1,i})\subseteq A_{n,i},
\qquad
\text{且在同一标号 \(i\) 上与 }\pi_{n+1\to n}\text{ 的作用共轭。}
\]
令 \(\Pi_{N\to n}:=\Pi_{n+1\to n}\circ\cdots\circ\Pi_{N\to N-1}\)（\(1\le n\le N\)）。
取 \(X_N,Y_N\) 为 \(A_N\) 上独立均匀随机元，并定义相干停时（带终止约定）
\[
\tau:=\min\{1\le n\le N:\ \Pi_{N\to n}(X_N)\neq \Pi_{N\to n}(Y_N)\},
\qquad
\text{若不存在则置 }\tau:=N+1.
\]
则对任意整数 \(0\le n\le N\) 有精确尾分布
\begin{equation}\label{eq:xi-coherence-stopping-time-tail}
\PP(\tau>n)=\frac{1}{k\,4^n}.
\end{equation}
从而点分布为
\[
\PP(\tau=1)=1-\frac{1}{4k},\qquad
\PP(\tau=n)=\frac{3}{k\,4^n}\ (2\le n\le N),\qquad
\PP(\tau=N+1)=\frac{1}{k\,4^N}.
\]
并且期望停时满足闭式
\begin{equation}\label{eq:xi-coherence-stopping-time-mean}
\EE[\tau]=1+\frac{1}{3k}\Bigl(1-4^{-N}\Bigr).
\end{equation}
\end{theorem}
\begin{proof}
事件 \(\{\tau>n\}\) 等价于 \(\Pi_{N\to n}(X_N)=\Pi_{N\to n}(Y_N)\)。
由 \(A_N\) 的 \(k\) 分量直和结构，该事件必先要求两次抽样落在同一分量（概率 \(1/k\)）。
在该分量内，由截断共轭性与均匀性，\(\Pi_{N\to n}(X_N)\) 在 \(A_{n,i}\) 上均匀；而 \(A_{n,i}\) 作为 \(G_n\)-挠空间与 \(G_n\) 等势，故两次独立均匀抽样相等的概率为 \(1/|G_n|=4^{-n}\)。
相乘得到 \eqref{eq:xi-coherence-stopping-time-tail}。
点分布由差分 \(\PP(\tau=n)=\PP(\tau>n-1)-\PP(\tau>n)\) 给出；期望由
\(
\EE[\tau]=\sum_{n\ge 0}\PP(\tau>n)=1+\sum_{n=1}^{N}(k4^n)^{-1}
\)
求和得到 \eqref{eq:xi-coherence-stopping-time-mean}。证毕。
\end{proof}

\begin{corollary}[Lee--Yang 五平行族的首层区分律与几何尾]\label{cor:xi-coherence-stopping-time-leyang-k5}
\leanverified{paper\_xi\_coherence\_stopping\_time\_leyang\_k5}
在推论 \ref{cor:xi-terminal-zm-leyang-ramification-pullback-2n} 的 Lee--Yang 分歧塔上，
将五个有限分歧点 \(P_0,\dots,P_4\) 的平移族
\(
P_i+E_{\min}[2^n]
\)
视为 \(k=5\) 个互不相交的 \(E_{\min}[2^n]\)-挠分量，并在各分量内用倍点截断诱导兼容的 \(4\)-叉截断塔。
则对任意固定深度 \(N\ge 1\)，相应相干停时 \(\tau\) 满足
\[
\PP(\tau=1)=\frac{19}{20},\qquad
\PP(\tau=n)=\frac{3}{5\cdot 4^{n}}\ (2\le n\le N),\qquad
\PP(\tau=N+1)=\frac{1}{5\cdot 4^{N}},
\]
并且
\[
\EE[\tau]=1+\frac{1}{15}\Bigl(1-4^{-N}\Bigr).
\]
\end{corollary}
\begin{proof}
将定理 \ref{thm:xi-coherence-stopping-time-k-parallel-4ary} 取 \(k=5\) 并用 \(|E_{\min}[2^n]|=4^n\) 代入即可。证毕。
\end{proof}

\begin{proposition}[分辨率读出与信息增益的严格恒等式]\label{prop:xi-coherence-resolution_information_identity}
\leanverified{paper\_xi\_coherence\_resolution\_information\_identity}
在定理 \ref{thm:xi-coherence-stopping-time-k-parallel-4ary} 的设定下，固定 \(N\ge 1\)。
在每个分量 \(A_{n,i}\) 上选择一次性标架同构 \(\theta_{n,i}:A_{n,i}\to G_n\)，并令坐标读出
\[
O_n:=\theta_{n,I_n}\!\bigl(\Pi_{N\to n}(X_N)\bigr)\in G_n,\qquad 1\le n\le N,
\]
其中 \(I_n\) 为 \(\Pi_{N\to n}(X_N)\) 所在分量标号。
则以 \(\log_2\) 为底的 Shannon 量度下，对任意 \(1\le n\le N\) 有
\[
I(X_N;O_n)=H(O_n)=\log_2|G_n|=2n,
\]
并且条件熵满足闭式
\[
H(X_N\mid O_n)=\log_2(k4^{N-n})=\log_2 k+2(N-n).
\]
\end{proposition}
\begin{proof}
由于 \(O_n\) 为 \(X_N\) 的确定函数，故 \(I(X_N;O_n)=H(O_n)\)。
由截断共轭性与分量均匀性，\(O_n\) 在 \(G_n\) 上均匀，故 \(H(O_n)=\log_2|G_n|=2n\)。
另一方面，给定 \(O_n=o\)，其原像由 \(k\) 个分量的均匀纤维组成，每个分量纤维大小为 \(|\ker(\pi_{N\to n})|=4^{N-n}\)，故条件分布为等概率于 \(k4^{N-n}\) 个点上，得到所示条件熵。证毕。
\end{proof}

\begin{theorem}[相干等价类的仿射子立方结构]\label{thm:xi-coherence-affine-subcube}
在 Lee--Yang 分歧塔的任一固定分量 \(P_i+E_{\min}[2^N]\) 内，选定 \(T_2(E_{\min})\) 的 \(\ZZ_2\)-基并据此得到同构
\(
E_{\min}[2^N]\cong(\ZZ/2^N\ZZ)^2
\)
与地址标识
\(
\Gamma_N:P_i+E_{\min}[2^N]\to\{0,1,2,3\}^N
\)
（将每一层的两位二进地址合并为一位 \(4\)-进数字）。
对 \(1\le m\le N\) 定义“分辨率 \(m\)”相干关系为
\[
x\sim_m y
\quad:\!\!\Longleftrightarrow\quad
\Gamma_N(x)\equiv \Gamma_N(y)\pmod{4^m},
\]
其中同余按 \(\sum_{t=0}^{N-1}\Gamma_N(\cdot)_t4^t\in\ZZ/4^N\ZZ\) 的自然解释给出。
则对每个 \(m\)：
\begin{enumerate}
  \item 等价类 \([x]_{\sim_m}\) 在 \(2N\)-维超立方体顶点模型中为一个维数 \(2(N-m)\) 的仿射子立方，故
  \[
  |[x]_{\sim_m}|=4^{N-m};
  \]
  \item 该分量被分割为 \(4^m\) 个两两不交的上述仿射子立方。
\end{enumerate}
\leanverified{paper\_xi\_coherence\_affine\_subcube}
\end{theorem}
\begin{proof}
在地址模型 \(\{0,1,2,3\}^N\) 中，同余 \(\bmod 4^m\) 等价于前 \(m\) 个 \(4\)-进位（即 \(t=0,\dots,m-1\) 的低位数字）一致，而余下 \(N-m\) 个数字可自由取值。
将每个 \(4\)-进位展开为两位二进坐标后，“固定若干坐标而其余坐标自由”即为仿射子立方的定义；
自由坐标数为 \(2(N-m)\)，顶点数为 \(2^{2(N-m)}=4^{N-m}\)，并且可能的固定低位共有 \(4^m\) 种。证毕。
\end{proof}

\begin{theorem}[Gödel \(4\)-进编码与相干停时的 \(4\)-进指数表述]\label{thm:xi-coherence-godel-4adic-valuation}
在定理 \ref{thm:xi-coherence-affine-subcube} 的地址标识下，定义 Gödel 型编码
\[
\mathbf G_N(x):=\sum_{t=0}^{N-1}\Gamma_N(x)_t\,4^t\in \ZZ/4^N\ZZ.
\]
则对任意 \(x,y\) 与任意 \(1\le m\le N\) 有严格等价
\[
x\sim_m y
\quad\Longleftrightarrow\quad
\mathbf G_N(x)\equiv \mathbf G_N(y)\pmod{4^m}.
\]
并且若定义停时
\[
\tau(x,y):=\min\{1\le m\le N:\ x\not\sim_m y\},\qquad
\text{若不存在则置 }\tau(x,y):=N+1,
\]
则
\[
\tau(x,y)=1+\nu_4\!\bigl(\mathbf G_N(x)-\mathbf G_N(y)\bigr),
\]
其中 \(\nu_4(t):=\max\{r\in\{0,1,\dots,N\}: t\equiv 0\ (\mathrm{mod}\ 4^r)\}\) 为截断到深度 \(N\) 的 \(4\)-进指数。
\leanverified{paper\_xi\_coherence\_godel\_4adic\_valuation}
\end{theorem}
\begin{proof}
由 \(\mathbf G_N\) 的定义，模 \(4^m\) 的同余恰等价于低位 \(m\) 个数字一致，即 \(\sim_m\) 的定义。
停时 \(\tau(x,y)\) 是使低位前缀首次失配的层数，等价于“最大相同低位前缀长度”，而该长度正由 \(\nu_4(\mathbf G_N(x)-\mathbf G_N(y))\) 给出。证毕。
\end{proof}

\endinput
